
\clearpage
% ============================================================================
\section*{v4.10.0-r75 Source-Only: Dated Register of the Model State and Its Predictions for Version v5.6.0 (Stand 28 November 2026)}
\addcontentsline{toc}{section}{v4.10.0-r75 Source-Only: Dated Register of the Model State and Its Predictions for Version v5.6.0 (Stand 28 November 2026)}
\label{sec:v4100-r75-prediction-register}
% ============================================================================

\paragraph{Purpose.}
On 28 November 2026 the author set a standing rule for every public version: each version records its model state together with its
predictions, with date and DOI, and older versions stay visible.  This layer states the rule (r75.1), lists the inputs of the model state of
version v5.6.0 with their status (r75.2), and records all open pre-registered predictions with the inputs each depends on, the data that decide
it and what would refute it (r75.3, r75.4).  It is a register, not a test: it adds no computation, changes no criterion and no verdict, and
every value is copied from the cited advance or result note, where the binding criterion stands.  Names follow r69: rest frame R (formerly
reference frame), matter frame M, dual frame T (formerly tachyon frame), light point P (formerly photon eigenframe); Re sector (neutral
leptons, \(c=0\)) and Im sector (charged leptons, \(c=1\)), under H99.  Earlier layers are not rewritten.

\noindent\textbf{Source-only status.}
This v4.10.0-r75 layer is source-only.  It does not build a PDF, does not assign a new DOI,
and does not authorize a public release.  The v5.5.0 Version DOI
\href{https://doi.org/10.5281/zenodo.24563811}{10.5281/zenodo.24563811} and the Concept DOI
\href{https://doi.org/10.5281/zenodo.19210728}{10.5281/zenodo.19210728} are unchanged by it.  The register is written for version v5.6.0
(Version DOI to be reserved); it becomes the reference for tests from the date of publication of v5.6.0.

\subsection*{v4.10.0-r75 the rule}
\label{subsec:v4100-r75-rule}

\paragraph{Definition v4.10.0-r75.1 (dated model state per version; standing rule of the author, 28 November 2026).}
(a) Every public version (Companion, physics article, German preprint) contains a dated section with the model state and its predictions:
date, DOI, and the list of predictions.  Older versions and their registers stay visible under their own DOIs; the register of v5.6.0 is the
first one, and versions v5.5.0 and earlier stay visible as published.
(b) For every prediction it is marked openly which inputs are derived, which are set (postulate or setting), which are measured inputs, and
what was found after the fact (post hoc).
(c) A correction is made only in a new version, with a stated reason: progress of the model or an error.  It is never a readjustment after a
new measurement.
(d) The version that was fixed before the measurement counts as the test.  Inside the repository the date of fixing is the commit of the
advance note; each note states which data were known and do not count.
(e) The register in the Companion is authoritative; the physics article and the German preprint refer to it.

\subsection*{v4.10.0-r75 model state}
\label{subsec:v4100-r75-model-state}

\paragraph{Register v4.10.0-r75.2 (inputs of the model state, Stand 28 November 2026).}
\begin{center}
\small
\begin{tabular}{@{}>{\raggedright\arraybackslash}p{0.30\textwidth}>{\raggedright\arraybackslash}p{0.17\textwidth}>{\raggedright\arraybackslash}p{0.45\textwidth}@{}}
\toprule
input & class & source and remark\\
\midrule
counting rule (weight \(2\) per spin, Casimir area, one bit for T) & postulate & adopted 26 November 2026 (r67); OP1 reads ``derive the counting postulate''\\
\(\beta_{\rm sym}=0.5494056410\), \(\gamma_L=1.8646997646\), \(\gamma_M=\gamma_{\rm BI}=0.2375330\) & computed from the postulate & r67, r68; not derived independently of the postulate\\
equal standing of the frames \(J\), hence \(K=\tfrac23\) & postulate & r63.7, adopted 26 November 2026; \(K=\tfrac23\) follows under \(J\)\\
Brannen form in both lepton sectors, H99 sector assignment & setting & as in r73; negative lightest neutrino root (H99)\\
common rotation angle \(\varphi=0.0395773\) & measured input & the one measured number of the lepton picture (r73, r74); left open by the author (r74.D)\\
sign of the branch-choice term & setting (author) & r73.6; the coefficient \(q\) of H167 is only this sign, \(q\to0^-\) (H236, r74.10)\\
scales \(\mu_e\), \(\mu_\nu\) & fitted (measured) & \(\mu_e\) from the charged pole masses, \(\mu_\nu\) from \(\Delta m^2_{31}\) (H227); ratio not a frame number\\
pole masses \(m_e,m_\mu,m_\tau\); \(\Delta m^2_{21}\), \(\Delta m^2_{31}\); \(m_Z\); mixing data for \(m_\beta\), \(m_{\beta\beta}\) & measured inputs & as cited in the source notes\\
placeholders CT-001, CT-002 (H102), CT-004 (H15, H46, H96); setting S1 (H125); setting V1 (H105) & settings / placeholders & guards of the counter-term register; not derived\\
\bottomrule
\end{tabular}
\end{center}
Circulant sector and tree level are assumed for every lepton entry.  The identification \(\Omega_M=\beta_{\rm sym}^2\) (H96) and the readings of
\(\beta_{\rm sym}^2\) as \(|U_{e2}|^2\) (H9) are additional identifications, not consequences of the postulate.

\subsection*{v4.10.0-r75 predictions}
\label{subsec:v4100-r75-predictions}

\paragraph{Register v4.10.0-r75.3 (open pre-registered predictions, values and tests).}
``Fixed'' is the date of the advance note (or, where stated, of the result note).  Experiment sensitivities are planning values of the
collaborations taken from the literature; they were not recomputed in the repository.

{\small
\begin{longtable}{@{}>{\raggedright\arraybackslash}p{0.14\textwidth}>{\raggedright\arraybackslash}p{0.20\textwidth}>{\raggedright\arraybackslash}p{0.15\textwidth}>{\raggedright\arraybackslash}p{0.20\textwidth}>{\raggedright\arraybackslash}p{0.21\textwidth}@{}}
\toprule
quantity & predicted value & source, fixed & deciding data & refuted if\\
\midrule
\endfirsthead
\toprule
quantity & predicted value & source, fixed & deciding data & refuted if\\
\midrule
\endhead
\bottomrule
\endfoot
mass ordering & normal (inverted branch \(-16\sigma\) from \(x\), H213) & H227, H176; 28 and 26 Nov 2026 & JUNO, DUNE, Hyper-K, cosmology & inverted ordering established at \(3\sigma\) (refutes H227 and H176)\\
\(\sum m_\nu\) & \(59.30\pm0.24\) meV & H227 (\texttt{massen\_\allowbreak bilanz\_\allowbreak h227}), 28 Nov 2026 & DESI DR3, Euclid, CMB-S4 (\(\sigma\approx15\)--\(20\) meV) & measured \(\sum m_\nu<58\) meV, or inverted ordering\\
\(m_1,m_2,m_3\) & \(0.379\pm0.002\), \(8.808\pm0.036\), \(50.11\pm0.20\) meV & H227, 28 Nov 2026 & indirect, through \(\sum m_\nu\) & as \(\sum m_\nu\)\\
\(m_\beta\) & \(8.87\pm0.04\) meV & H227, 28 Nov 2026 & KATRIN (\(\approx300\) meV), Project 8 (goal \(\approx40\) meV) & any detection at these sensitivities (null result predicted)\\
\(m_{\beta\beta}\) & \(1.29\) to \(4.00\) meV (Majorana phases free); \(3.49\) meV in the Brannen sign reading & H227, 28 Nov 2026 & LEGEND-1000, nEXO (\(\approx9\)--\(20\) meV) & a signal above \(\approx10\) meV (null result predicted)\\
\(R\) \(=\Delta m^2_{21}/\Delta m^2_{31}\) & A (H99) \(0.030839\); B \(0.030087\); C \(0.030375\) & H100, 13 Nov 2026 (H227 repeats A, not blind) & first NuFIT version after 6.1, NO, with SK; void without one by end of 2027 & per candidate \(|z|>2\); H100 confirmed if A refuted and B or C compatible\\
\(\sin^2\theta_{12}\), reading H9 & \(\beta_{\rm sym}^2/\cos^2\theta_{13}\) \(=0.30862\pm0.00018\) & H9, 20 Oct 2026 & next JUNO result with \(\sigma\le0.004\) & \(|z|\ge3\)\\
\(\sin^2\theta_{12}\), reading H176 & \(0.3019\) (band \(0.3012\)--\(0.3025\)) & H176, 26 Nov 2026 & JUNO, NuFIT & deviation \(>3\sigma\) (measurement and band)\\
\(\sin^2\theta_{12}\), structural reading H206--H208 & \(\approx0.318\)--\(0.3195\) (near TM1); \(0.3155\) only with an allowed, non-derived extra term and \(|\sin\delta|\approx1\) & H208 result note, 27 Nov 2026 & JUNO with DUNE (correlation \(\theta_{12}\leftrightarrow\delta_{\rm CP}\)) & JUNO central near \(0.3036\) together with maximal CP violation\\
\(\sin^2\theta_{23}\) & \(\le0.54\) (H42); \([0.49,0.50]\) (H46); \(0.4954\) (H176); \(0.49445\) (H207) & H42, H46 24 Oct 2026; H176 26 Nov; H207 result note 27 Nov & first global fit or DUNE/Hyper-K result with a \(2\sigma\) interval & H42: \(2\sigma\) interval above \(0.54\); H46: \(2\sigma\) interval misses \([0.49,0.50]\); H176: \(>3\sigma\)\\
\(\sin^2\theta_{13}\) & \(0.0223\) (H176, band \(0.0218\)--\(0.0228\)); \(0.02164\) without parameter (H207) & H176 26 Nov; H207 result note 27 Nov & reactors, global fits & H176: \(>3\sigma\)\\
\(\delta_{\rm CP}\) & \(\pm77.7^\circ\), \(|\sin\delta|=0.977\) (H176); \(\sin\delta<0\) (H125); together \(\approx282^\circ\) & H176 26 Nov; H125 24 Nov 2026 & DUNE, Hyper-K, global fits; H125 until 31 Dec 2035 & H176: \(|\sin\delta|<0.8\) at \(3\sigma\); H125: \(\sin\delta\le0\) excluded at \(\ge3\sigma\)\\
\(\Omega_m\) (flat \(\Lambda\)CDM) & \(\beta_{\rm sym}^2=0.30185\) & H96, 6 Nov 2026 & first DESI DR3 (BAO alone) or Euclid value after 6 Nov 2026 & deviation \(>2\sigma\)\\
\(\sin^2\theta_{12}\cos^2\theta_{13}\) \(=\Omega_m\) & equality of the two places, whatever the value & H112, 17 Nov 2026 & the test data of H9 and H96, both; void without both by 31 Dec 2029 & \(|z|\ge3\) with the combined error (weak test)\\
\(m_W\) & \(80.3615\pm0.0018\) GeV & H15, 21 Oct 2026 & first world average or official combination after 21 Oct 2026 & \(|z|\ge3\) (tension for \(2\le|z|<3\))\\
\(m_\tau\) & \(1777.134\) MeV (H102) against exact Koide \(1776.969\) MeV (H227) & H102, 13 Nov 2026; H227 28 Nov 2026 (not blind) & first PDG edition after 13 Nov 2026 (until end of 2028); Belle II, BESIII & H102: \(|z|>2\); separates from Koide only at \(\sigma\lesssim0.05\) MeV\\
area steps and echoes (under V1) & \(\Delta A=4\ln2\,\ell_P^2\), absorption lines \(\omega_n=n\tfrac{\ln2}{2\pi}\kappa\) \(+\,m\Omega_H\) & H105, 16 Nov 2026 & LVK, later LISA analyses after 16 Nov 2026; until 31 Dec 2028 & excluded for one event (Bayes factor \(\le1/10\) or \(95\%\))\\
\end{longtable}
}

\paragraph{Register v4.10.0-r75.4 (what each prediction depends on).}
\emph{Derived} means derived inside the model from the stated settings; \emph{set} means postulate, setting or fitted choice; \emph{measured}
means a measured input; \emph{post hoc} marks what was found after the data were known.

{\small
\begin{longtable}{@{}>{\raggedright\arraybackslash}p{0.12\textwidth}>{\raggedright\arraybackslash}p{0.20\textwidth}>{\raggedright\arraybackslash}p{0.24\textwidth}>{\raggedright\arraybackslash}p{0.17\textwidth}>{\raggedright\arraybackslash}p{0.17\textwidth}@{}}
\toprule
entry & derived & set & measured & post hoc / blindness\\
\midrule
\endfirsthead
\toprule
entry & derived & set & measured & post hoc / blindness\\
\midrule
\endhead
\bottomrule
\endfoot
H227 (ordering, \(\sum m_\nu\), \(m_i\), \(m_\beta\), \(m_{\beta\beta}\)) & spectrum from the Brannen form at \(\pi/6-\varphi\); offset as a common rotation (r73.4); \(m_\beta\), \(m_{\beta\beta}\) range & \(K=\tfrac23\) (postulate \(J\)); Brannen form, H99; branch sign (author); \(q\to0^-\) (H236); scale \(\mu_\nu\) fitted; sign reading for \(3.49\) meV & \(\varphi\); \(\Delta m^2_{31}\); mixing data & masses not blind (H213, H168); \(m_\beta\), \(m_{\beta\beta}\) new; DESI DR2 tendency noted post hoc\\
H99/H100 (\(R\)) & \(R\) from \(\theta_\nu=\theta_e+\pi/12\) (A) without free number & H99 (A); one scale (B); Sumino \(c=1\) (C) & pole masses; NuFIT 6.1 for B, C & A not blind against NuFIT 6.0/6.1 (\(+1.31\sigma\), \(+1.80\sigma\)); B, C chosen after NuFIT 6.1\\
H9 (\(\theta_{12}\)) & --- & reading \(|U_{e2}|^2=\beta_{\rm sym}^2\); \(\beta_{\rm sym}\) from the counting postulate & \(\theta_{13}\) & relation found post hoc; printed value contains the reference-value result \(0.30868\) (r68.2)\\
H176 (\(\theta_{12}\), \(\theta_{13}\), \(\theta_{23}\), \(\delta\), ordering) & mixing matrix from the fixed candidate & one-field form, \(K=\tfrac23\); direction \((K_1+K_2)/\sqrt2\) and \(x=0.04478\) fitted & pole masses; \(\Delta m^2_{21}\) & selected post hoc from eight variants and fitted to today's data (\(\chi^2=2.86\)); counts only against future data\\
H206--H208 (\(\theta_{12}\leftrightarrow\delta\)) & fixed column \(v_0\), \(|U_{e1}|^2=\tfrac23\) exact (Wolfram); direction \(K_2\) (H206); TM1 tree (H173) & one-field form; extra term S+d allowed, not derived & \(\theta_{13}\) (pinned) & correlation found post hoc (H208); result notes, no layer yet\\
H207 (\(\theta_{13}\), \(\theta_{23}\), \(\delta\)) & values with strength \(\tan\varphi\) in direction \(K_2\) & as H206--H208 & \(\varphi\) & partly not blind; result note, not a separate advance note for future data\\
H42, H46 (\(\theta_{23}\)) & bound and window from r55 with mask r57 & TM2 reading r55; mask \(W_l\); forms S/K; H46: \(\varepsilon=\operatorname{artanh}\beta_{\rm sym}\) (counting postulate); CT-004 (H46) & data of r57 windows & H42 bound noticed after evaluation; reading \(\varepsilon=\eta\) not chosen blind\\
H125 (\(\operatorname{sign}\sin\delta\)) & relative sign of leptons and quarks (stage A, exact) & setting S1 (CP from the family structure in both sectors) & sign of \(J_{\rm CKM}\) & stage A not blind; a single sign has prior \(\tfrac12\)\\
H96 (\(\Omega_m\)) & \(\beta_{\rm sym}^2\) from the postulate & identification \(\Omega_M=\beta^2\); epoch ``today''; CT-004 & --- & no free parameter; known DESI/Planck values do not count\\
H112 & --- & the two identifications of H9 and H96 & --- & independent of the value of the counting rule\\
H15 (\(m_W\)) & \(m_W=m_Z\gamma_L/\sqrt{1+\gamma_L^2}\) given the relation & relation \(\cot\theta_W=\gamma_L\); \(\gamma_{\rm BI}\) from the counting equation (postulate); CT-004 & \(m_Z\) & relation found post hoc (\(0.8\sigma\) at the known value)\\
H102 (\(m_\tau\)) & \(m_\tau\) from \(\Phi_e=\pi/4-\gamma_M/2\) & relation; \(\gamma_M\) from the counting rule; CT-001, CT-002 & \(m_e\), \(m_\mu\) & found post hoc in a scan (154 trials); today \(+2.3\sigma\)\\
Koide \(m_\tau\) (H227) & \(m_\tau\) from \(m_e,m_\mu\) at \(K=\tfrac23\) & postulate \(J\) & \(m_e\), \(m_\mu\) & not blind (\(+0.43\sigma\))\\
H105 & lines and steps under the crossing reading & setting V1 (each horizon crossing a T quantum), chosen to give a measurable consequence & --- & \(\alpha=4\ln2\) not specific to the project\\
\end{longtable}
}

\paragraph{Remark v4.10.0-r75.5 (scope of the register).}
(a) The tests H173--H212 behind several entries (H176, H206--H208, H213) are documented in research notes and are not yet recorded as a layer,
as r73 states; entering them in this register makes no claim beyond their result notes.  (b) H227 is recorded in r73.10, H236 in r74.10.
(c) A prediction that is withdrawn or corrected in a later version stays listed here as fixed; the later version states the reason.
(d) The rows do not form independent tests: H15, H46 and H96 rest on the same counting core (v4.10.0-r08, r10); H9, H96 and H112 share
\(\beta_{\rm sym}^2\); the three readings of \(\theta_{12}\) exclude one another, so JUNO decides between them; the neutrino rows of H227 and the
ordering criterion of H176 fall together with an established inverted ordering.  (e) Data values known when a note was fixed are listed in
that note and do not count as tests.

\subsection*{v4.10.0-r75 frames}
\label{subsec:v4100-r75-frames}

\begin{center}
\small
\begin{tabular}{@{}>{\raggedright\arraybackslash}p{0.34\textwidth}>{\raggedright\arraybackslash}p{0.22\textwidth}>{\raggedright\arraybackslash}p{0.36\textwidth}@{}}
\toprule
item & frame & comparison\\
\midrule
lepton masses, mixing, \(R\), \(m_\tau\), \(m_W\) & M & compared with data in M, as in the source notes\\
H42, H46 (mask with \(\varepsilon\)) & R \(\to\) M through the declared mask & as in r57, r60\\
\(\Omega_m\) (H96, H112) & M (cosmological readout) & identification stated in H96\\
H105 & horizon crossing (T quanta under V1) & as in r24.4\\
P & not used & ---\\
\bottomrule
\end{tabular}
\end{center}
No new frame map is used; every entry keeps the frame and the identification of the note that fixed it.

\subsection*{v4.10.0-r75 register entries}
\label{subsec:v4100-r75-register}

\begin{center}
\small
\begin{tabular}{@{}>{\raggedright\arraybackslash}p{0.40\textwidth}>{\raggedright\arraybackslash}p{0.52\textwidth}@{}}
\toprule
statement & status after v4.10.0-r75\\
\midrule
rule ``dated model state per version'' (r75.1) & standing rule of the author, 28 November 2026\\
register r75.2--r75.4 & reference for tests from the publication of v5.6.0 (Version DOI to be reserved)\\
H9, H15, H42, H46, H96, H100, H102, H105, H112, H125, H176, H227; readings H206--H208 & open predictions; criteria unchanged\\
versions v5.5.0 and earlier & stay visible under their DOIs\\
OP1--OP5 & open\\
\bottomrule
\end{tabular}
\end{center}

\subsection*{v4.10.0-r75 notes}
\label{subsec:v4100-r75-notes}

{\raggedright
Sources: \texttt{wiki/\allowbreak Pruefbare-Vorhersagen.md}, table ``Offene Vorhersagen'' in \texttt{wiki/\allowbreak Vorhersagen-und-Tests.md};
advance notes in \texttt{notes/\allowbreak research/\allowbreak predictions/} (H42, H46, H96, H99, H100, H102, H105, H112, H125),
\texttt{notes/\allowbreak research/\allowbreak v4\_\allowbreak 8\_\allowbreak 0\_\allowbreak r03\_\allowbreak r12/\allowbreak llm/} (H9),
\texttt{notes/\allowbreak research/\allowbreak v4\_\allowbreak 9\_\allowbreak 0\_\allowbreak r01\_\allowbreak r16/\allowbreak llm/} (H15);
\texttt{notes/\allowbreak research/\allowbreak hilbert\_\allowbreak shift\_\allowbreak felder\_\allowbreak 2026\_\allowbreak 11\_\allowbreak 24/}: folders
\texttt{mischung\_\allowbreak h176/}, \texttt{winkel\_\allowbreak in\_\allowbreak ebene\_\allowbreak h206/}, \texttt{staerke\_\allowbreak h168\_\allowbreak h207/},
\texttt{theta12\_\allowbreak tm1\_\allowbreak h208/}, \texttt{versatz\_\allowbreak spannung\_\allowbreak h213/}, \texttt{massen\_\allowbreak bilanz\_\allowbreak h227/},
\texttt{doppelzaehlung\_\allowbreak h236/}.  Every value was checked against the cited note.
Layer notes: \texttt{notes/\allowbreak research/\allowbreak v4\_\allowbreak 10\_\allowbreak 0\_\allowbreak r75\_\allowbreak vorhersage\_\allowbreak register/}.
No computation; no Qiskit, no HDL/MMA comparison.
\par}

\subsection*{v4.10.0-r75 ledger and guard update}
\label{subsec:v4100-r75-ledger-guard}

\[
R_{\rm chk}^{v4100\text{-}r75}=(1,0,0)
\]
for (status recorded; no computation; no comparison with data).

The layer spends no guard:
\[
R_{\rm glob}^{(v4100\text{-}r75,+)}=(0,1,1,1,1),
\qquad
\Gamma_{\rm guard}^{v4100\text{-}r75}=(1,1,1,1,1),
\qquad
\beta_4=1/256\ \text{unspent}.
\]
The layer does not imply exact finite \(C5\) zero closure, exact all-readout terminality, \(C6/T2\), a global photon-clock theorem,
or global all-readout closure.  OP1--OP5, C5 and C6/T2 remain open.
